\section{Main Concepts}

% \begin{itemize}
%   \item GMaps
%   \item Graph transformation rules
%   \item Orbit tracking
% \end{itemize}

% In this section, we present the generalized maps, graph transformation rules and
% their subsequent concepts. Which are necessary to the understanding of our
% contribution.

\subsection{Generalized maps}

\emph{Generalized maps} (or G-maps) \cite{lienhardt1991topological,
  Damiand-Lienhardt14} allow the representation of manifold geometric objects
(with or without boundaries), based on a cellular $n$-dimensional topological
structure. %
The representation of an object as a G-map comes intuitively from its
decomposition into topological cells (vertices, edges, faces, volumes, and so
on). The lowest topological entity in a G-map is called a dart. In a G-map, such
a topological cell is represented as a subgraph called an orbit. %

\subsection{Graph transformation rules}

% Jerboas's \cite{belhaouari2014jerboa} graph transformation rules allow the
% formalization of operations over G-maps. %

% In a few words, a \emph{rule} $r : L \longrightarrow R$ and a \emph{match}
% $m : L \to G$ to a G-map $G$, describe the transformation
% $G \longrightarrow^{r,m} H$ from $G$ to $H$. %
% The match $m$ allows the replacement of a subgraph of $G$ described by the
% left-hand side of the rule $L$ with another one described by the right-hand side
% $R$, in order to produce $H$. %

\subsection{Orbit tracking}

% Rules contain the necessary information characterizing the topological changes
% affecting an orbit throughout an application. %
% Thus, the tracking of an orbit is automatically made without any addition other
% than the rules' syntactic analysis. %

\section{Parametric Specification }

\begin{itemize}
  \item persistent name
%   \item parametric specification syntax
\end{itemize}

% During an object's construction, a \emph{parametric specification} records both
% the rules representing the operations applied and their parameters (both
% topological and geometric) in order to describe the modeling process. %
% Editing a parametric specification means that rules may be added, deleted,
% moved, and their parameters can be changed. %
% With such changes, the topological parameters may have different concrete names,
% be deleted and so on. %
% As a result, using the parameters' concrete names eventually lead to unexpected
% results (at best) or failure at runtime. %


\subsection{Persistent name}

While Jerboa's graph transformation rules use darts as topological parameters,
it follows that each persistent name must represent a unique dart. As it
happens, rules make it possible to determine unambiguously which node filters or
creates any dart.

% Since rules use darts as topological parameters, it follows that each persistent
% name must represent a unique dart. %
% As it happens, rules make it possible to determine unambiguously, which node
% filters or creates any dart. %


% \subsection{Parametric Specification syntax}

% A number of fields are used to describe an application within a parametric
% specification, namely an application number, a rule application with
% the topological parameters and the geometric parameters. %

\section{Evaluation }

% Although a persistent name represents the history of a dart, an orbit is subject
% to topological changes and, thus, requires the construction of its own
% history in order to be accurately matched at reevaluation. %
Once the initial evaluation is done and its parametric specification has been
built, an evaluation's Directed Acyclic Graph (or DAG) must be issued for each
persistent name before any parametric specification can be reevaluated. %
An evaluation DAG traces the history of each topological parameter back to the
first created orbits it originates from, thus allowing the matching of the
corresponding topological parameter at reevaluation time. %


\subsection{Evaluation DAG}

An evaluation DAG is built by parsing the applications and nodes of a persistent
name from end to start. It is sorted by levels representing the different
applications inside a history. Each level is made of an orbit level and an event
level. The orbit level contains a node’s name and some orbits. The event level
contains a rule’s application number and some events.

\begin{figure}[H]
  \centering
  \includegraphics[width=0.3\textwidth]{pn7}
  \caption{$PN_{7}$'s evaluation DAG}
  \label{fig:eval-dag-pn7}
\end{figure}


\subsection{Traces and origins}

In order to accurately represent the history of an orbit, two types of arrows
are used in an evaluation DAG (and later in the reevaluation DAG). A black arrow
or trace arrows represents the evolution of an orbit throughout the modeling
process. A red arrow, or origin arrow, represents the origin of traced orbit in
the form of another orbit. In Fig.~\ref{fig:eval-dag-pn7}, the $\orbit{0,1}$ at
orbit level 2 is related to two orbits at orbit level 2, the $\orbit{0,1}$ orbit
with a trace arrow and the $\orbit{0}$ with an origin arrow. The trace arrow
relation here means the orbit $\orbit{0,1}$ at orbit level 3 is an evolution of
the orbit $\orbit{0,1}$ at orbit level 2. The origin arrow is added to reflected
that since the orbit $\orbit{0,1}$ is split by the \texttt{3-triangulation} rule application, the orbit tracked in the DAG comes from the orbit $\orbit{0}$ at orbit level 2.


\section{Reevaluation }

Each evaluation DAG represents an orbit's history which is valid with regards to
the initial evaluation. %
When reevaluating, editing the parametric specification makes the topological
parameters subject to changes. %
Thus, it is necessary to build reevaluation DAGs from the evaluation DAGs in
order to update topological parameters. %
Once built, a reevaluation DAG can designate one, several, or no orbit depending
on the editing of the parametric specification.

\subsection{Reevaluation DAG}

Contrary to the related evaluation DAG, a reevaluation DAG is built
top-down throughout the reevaluation process. While an evaluation
DAG represents the orbit’s history of a topological parameter, the
reevaluation DAG represents the history of this very same orbit af-
ter editing the parametric specification.

\begin{figure}[H]
    \centering
    \includegraphics[width=\columnwidth]{pn7-reval}
    \caption{$PN_7$ reevaluation DAG}
    \label{fig:pn7-reval-dag}
\end{figure}

\subsection{Parameter matching strategies}

The editing of the parametric specification leads to having a different DAG at
reevaluation (with event levels and/or branches being added, deleted or both).
For example, an orbit split present in the evaluation DAG may disappear during
the reevaluation, a merging can be added and so on. Several matching strategies
can then be considered depending on the application’s context such as in Fig.~\ref{fig:pn7-reval-dag}.


% %%% Local Variables:
% %%% mode: latex
% %%% TeX-master: "../main"
% %%% End:
